% ============================================================================
% Chapter 3: Cardano Improvement Proposals (CIPs)
% Study Guide — Blockchain-Based Credential Verification
% ============================================================================

\chapter{Cardano Improvement Proposals (CIPs)}
\label{ch:cips}

% --- Shared TikZ styles for this chapter ------------------------------------
\tikzset{
  guidebox/.style={draw, rounded corners=4pt, thick, align=center,
    font=\sffamily\small, fill=white, minimum height=1cm},
  guidearrow/.style={-{Latex[length=2.5mm]}, thick, color=gray!70!black},
  guidelabel/.style={font=\sffamily\footnotesize, text=gray!50!black},
}

Cardano's technical evolution is governed by a structured, community-driven
process called the \emph{Cardano Improvement Proposal} (CIP) system.  Every
significant feature---from how wallets derive keys to how smart contracts
reference on-chain data---begins life as a numbered CIP, undergoes public
review, and is either accepted, modified, or rejected before implementation.
Understanding the CIP landscape is essential for anyone building on Cardano
because CIPs define the \emph{interface contracts} that all tools, wallets,
and dApps must honour.

This chapter examines the eight CIPs most relevant to the credential
verification system described in this guide.  We begin with the governance
process itself, then walk through the token metadata standards (CIP-25 and
CIP-68), the wallet connectivity bridge (CIP-30), the Plutus~V2 trio that
made advanced datum handling possible (CIPs~31, 32, and~33), and the
hierarchical deterministic wallet derivation standard (CIP-1852).  For each
CIP we explain \emph{what problem it solves}, \emph{how it works
technically}, and \emph{how the credential system uses it}.


% ============================================================================
\section{What Are CIPs?}
\label{sec:cips-overview}
% ============================================================================

A Cardano Improvement Proposal is a formal design document that describes a
new feature, process, or informational guideline for the Cardano ecosystem.
The CIP process is modelled directly on the Internet Engineering Task Force's
\emph{Request for Comments} (RFC) system and Python's PEP process: anyone may
author a proposal, but acceptance requires community review, editor approval,
and---for protocol-level changes---implementation by the core development
teams at Input Output Global (IOG), the Cardano Foundation, or Emurgo.

CIPs are classified into three categories:

\begin{description}[leftmargin=1.5em, style=nextline]
  \item[Standards Track] Proposals that change the protocol, transaction
    format, validation rules, or networking layer.  Examples include CIP-31
    (Reference Inputs) and CIP-32 (Inline Datums).
  \item[Process] Proposals that change governance procedures, the CIP process
    itself, or tooling requirements.
  \item[Informational] Proposals that document conventions, best practices, or
    design rationale without mandating changes.  CIP-25 (NFT Metadata) began
    as an informational standard that the community adopted voluntarily.
\end{description}

The lifecycle of a CIP proceeds through several stages: \textbf{Draft} (initial
submission to the CIP repository on GitHub), \textbf{Proposed} (editors
confirm the proposal is well-formed), \textbf{Active} (the community has
reached rough consensus and implementations exist), and optionally
\textbf{Inactive} or \textbf{Superseded} if a newer CIP replaces it.

For the credential system, CIPs are not merely background reading---they are
the \emph{specifications} that the code implements.  The minting policy encodes
CIP-25 metadata structures.  The wallet module follows CIP-1852 derivation
paths.  The revocation mechanism depends on CIP-68 datum semantics.
Understanding each CIP at a technical level is prerequisite to understanding
why the system works the way it does.

\begin{figure}[ht]
\centering
\begin{tikzpicture}[node distance=1.8cm and 0.5cm]
  \node[guidebox, fill=blue!8, minimum width=2.6cm] (draft) {Draft};
  \node[guidebox, fill=yellow!8, minimum width=2.6cm, right=of draft] (proposed) {Proposed};
  \node[guidebox, fill=green!8, minimum width=2.6cm, right=of proposed] (active) {Active};
  \node[guidebox, fill=gray!10, minimum width=2.6cm, below=0.9cm of active] (inactive) {Inactive /\\Superseded};

  \draw[guidearrow] (draft) -- node[above, guidelabel] {editor review} (proposed);
  \draw[guidearrow] (proposed) -- node[above, guidelabel] {consensus} (active);
  \draw[guidearrow] (active) -- node[right, guidelabel] {replaced} (inactive);
  \draw[guidearrow, dashed] (proposed.south) -- ++(0,-0.5) -| node[below, near start, guidelabel] {withdrawn} (draft.south);
\end{tikzpicture}
\caption{CIP lifecycle stages.  A proposal moves from Draft through Proposed
to Active once community consensus is reached and implementations exist.
Proposals may be withdrawn or eventually superseded by newer standards.}
\label{fig:cip-lifecycle}
\end{figure}


% ============================================================================
\section{CIP-25: NFT Metadata Standard}
\label{sec:cip25}
% ============================================================================

\subsection{The Problem}

Cardano's native multi-asset framework allows anyone to mint custom tokens,
but the ledger itself stores only three pieces of information about each
token: the \emph{policy ID} (a Blake2b-224 hash of the minting policy script),
the \emph{asset name} (an arbitrary byte string up to 32~bytes), and the
\emph{quantity}.  There is no built-in mechanism for attaching human-readable
names, images, descriptions, or any other metadata to a token.  Without a
standard, every marketplace, wallet, and explorer would need to negotiate
metadata formats bilaterally---an untenable situation for an open ecosystem.

\subsection{How It Works}

CIP-25 solves this by defining a JSON schema that is embedded in the
\textbf{transaction auxiliary data} (formerly called ``transaction metadata'')
at the time the token is minted.  The auxiliary data is a map keyed by
integer labels; CIP-25 reserves label~\texttt{721} for NFT metadata.

The schema is hierarchical: under label~721, the top-level key is the
\textbf{policy ID} (hex-encoded).  Under each policy ID, keys are
\textbf{asset names} (hex or UTF-8).  Under each asset name, the metadata
object contains at minimum:

\begin{itemize}[nosep]
  \item \texttt{name} --- a human-readable display name
  \item \texttt{image} --- a URI (typically IPFS) pointing to the token's
    image
  \item \texttt{description} --- free-text description
\end{itemize}

Additional fields are permitted and commonly used: \texttt{mediaType},
\texttt{files} (for multi-asset bundles), and arbitrary key-value pairs.
Wallets and marketplaces parse the minting transaction's auxiliary data to
display this information to users.

A critical property of CIP-25 metadata is that it is \textbf{immutable once
written}.  The metadata exists in the auxiliary data of a specific
transaction.  Transactions, once confirmed on-chain, cannot be altered.  This
means that whatever metadata is attached at mint time is the metadata
forever---there is no mechanism within CIP-25 to update, correct, or revoke
the metadata after minting.

For simple NFTs---artwork, collectibles, event tickets---immutability is a
feature.  For credentials that may need to reflect status changes
(revocation, suspension, renewal), immutability is a fundamental limitation.
This limitation is precisely what motivated CIP-68.

\subsection{How the Credential System Uses CIP-25}

In the credential system, CIP-25 metadata is attached to license NFTs at
mint time.  The metadata records the license type, jurisdiction, issue date,
license number, and the licensee's public key hash.  Because this
information is immutable, it serves as a permanent attestation: ``At this
point in time, this authority issued this credential to this individual.''

Verification of the attestation is straightforward: a verifier checks that
(1)~the token exists in the holder's wallet, (2)~the policy ID matches the
known authority, and (3)~the CIP-25 metadata in the minting transaction
matches the claimed credential details.

\begin{figure}[ht]
\centering
\begin{tikzpicture}[
  metabox/.style={guidebox, text width=5.5cm, align=left,
    font=\ttfamily\scriptsize, minimum height=0.6cm},
]
  % Outer transaction box
  \node[guidebox, fill=blue!8, minimum width=8.5cm, minimum height=8.2cm,
    label={[font=\sffamily\small\bfseries]above:Minting Transaction Auxiliary Data}] (txbox) at (0,0) {};

  % Label 721
  \node[guidebox, fill=yellow!8, minimum width=7.5cm, minimum height=7cm,
    anchor=north] (l721) at (0,3.3) {};
  \node[guidelabel, anchor=north west] at (l721.north west) {label: \textbf{721}};

  % Policy ID
  \node[guidebox, fill=green!8, minimum width=6.5cm, minimum height=5.8cm,
    anchor=north] (pid) at (0,2.5) {};
  \node[guidelabel, anchor=north west, text width=6cm] at (pid.north west)
    {policy\_id: \texttt{a1b2c3...}};

  % Asset name
  \node[guidebox, fill=orange!8, minimum width=5.5cm, minimum height=4.4cm,
    anchor=north] (aname) at (0,1.6) {};
  \node[guidelabel, anchor=north west] at (aname.north west)
    {asset\_name: \texttt{PE\_License\_042}};

  % Fields
  \node[metabox, fill=white] (f1) at (0,0.4)
    {\textcolor{keyword}{"name"}: "PE License \#042"};
  \node[metabox, fill=white, below=0.15cm of f1] (f2)
    {\textcolor{keyword}{"image"}: "ipfs://Qm..."};
  \node[metabox, fill=white, below=0.15cm of f2] (f3)
    {\textcolor{keyword}{"description"}: "California PE"};
  \node[metabox, fill=white, below=0.15cm of f3] (f4)
    {\textcolor{keyword}{"license\_type"}: "PE"};
  \node[metabox, fill=white, below=0.15cm of f4] (f5)
    {\textcolor{keyword}{"jurisdiction"}: "US-CA"};
  \node[metabox, fill=white, below=0.15cm of f5] (f6)
    {\textcolor{keyword}{"issue\_date"}: "2026-01-15"};

  % Immutable label
  \node[font=\sffamily\footnotesize\bfseries, text=red!60!black,
    rotate=35] at (3.2,2.6) {IMMUTABLE};
\end{tikzpicture}
\caption{CIP-25 metadata structure.  The JSON schema is nested under
label~721, then by policy ID and asset name.  Once written to the
transaction auxiliary data, the metadata cannot be changed.}
\label{fig:cip25-metadata}
\end{figure}


% ============================================================================
\section{CIP-68: Datum Metadata Standard}
\label{sec:cip68}
% ============================================================================

CIP-68 is the single most important CIP for the credential verification
system.  It solves the fundamental limitation of CIP-25---immutable
metadata---by introducing a \textbf{dual-token pattern} that separates
ownership from metadata, enabling authority-controlled updates without
requiring access to the token holder's private key.  This section is
intentionally the longest in the chapter because understanding CIP-68 is
prerequisite to understanding how credential revocation works.

\subsection{The Problem CIP-68 Solves}

Consider a professional engineering license issued as a CIP-25 NFT.  The
metadata at mint time includes \texttt{status: active}.  Two years later, the
licensing board discovers misconduct and needs to revoke the license.  Under
CIP-25, this is \emph{impossible on-chain}: the metadata is baked into the
minting transaction and cannot be altered.  The board's only options are:

\begin{enumerate}[nosep]
  \item \textbf{Burn and re-mint.}  The authority could burn the original
    token and mint a new one with \texttt{status: revoked}.  But burning
    requires consuming the UTxO that holds the token---which is in the
    \emph{licensee's} wallet.  Under the eUTxO model, spending a UTxO
    requires the holder's private key.  A revoked licensee has no incentive
    to cooperate, and the authority \emph{cannot} spend the licensee's UTxO
    without that key.
  \item \textbf{Off-chain revocation list.}  The authority maintains a
    database of revoked credentials, and verifiers must check both the
    on-chain token and the off-chain list.  This reintroduces centralized
    infrastructure, single points of failure, and availability dependencies
    ---exactly the problems blockchain was supposed to solve.
  \item \textbf{Accept that on-chain revocation is impossible.}  Many
    early Cardano credential projects reached this conclusion and treated
    the blockchain as a timestamping service only.
\end{enumerate}

CIP-68 provides a fourth option: \textbf{mutable on-chain metadata} through
a clever separation of concerns.

\subsection{The Dual-Token Pattern}

CIP-68 defines a naming convention using \textbf{asset name labels}---a
4-byte prefix prepended to the asset name that identifies the token's role.
The two most important labels are:

\begin{description}[leftmargin=2em, style=nextline, nosep]
  \item[Label 100 --- Reference Token]
    Held by the \textbf{issuing authority}.  This token sits at a script
    address (not a personal wallet) and carries an \textbf{inline datum}
    containing all mutable metadata: license type, jurisdiction, issue date,
    status, and any other fields the authority may need to update.  Because
    the Reference Token is at a script address controlled by the authority,
    the authority can spend and re-create this UTxO with an updated datum at
    any time, without any involvement from the token holder.

  \item[Label 222 --- User Token]
    Held by the \textbf{credential holder} (the licensee).  This token is
    a standard native asset that lives in the holder's personal wallet.  It
    is \emph{immutable} in the CIP-25 sense---the holder possesses it as
    proof of credential issuance.  The User Token contains no mutable
    metadata; its sole purpose is to attest that the holder was issued a
    credential under a specific policy.
\end{description}

The two tokens share the same policy ID and base asset name, differing only
in their 4-byte label prefix.  This creates a cryptographic link: given a
User Token with asset name \texttt{000de140} + \texttt{PE\_License\_042},
anyone can deterministically compute the corresponding Reference Token's
asset name as \texttt{000643b0} + \texttt{PE\_License\_042} and look up its
inline datum to read the current metadata.

\subsection{How Revocation Works}

The breakthrough insight of CIP-68 for credential systems is this: the
authority can \textbf{revoke a credential by updating the Reference Token's
datum} from \texttt{status: active} to \texttt{status: revoked}.  This
operation:

\begin{itemize}[nosep]
  \item Requires \textbf{only the authority's signature} (the Reference Token
    is at the authority's script address).
  \item Does \textbf{not touch the User Token} in the licensee's wallet---the
    authority never needs the licensee's private key.
  \item Propagates \textbf{globally within minutes} (Cardano's finality
    window).
  \item Is \textbf{verifiable by anyone}: read the Reference Token's inline
    datum, check the status field.
  \item Is \textbf{auditable}: the chain of datum updates (active $\to$
    suspended $\to$ revoked) is visible in the transaction history.
\end{itemize}

This is the mechanism that makes eUTxO-based credential revocation possible
without off-chain infrastructure, without the holder's cooperation, and
without burning the original token.

\subsection{Technical Details}

The Reference Token's inline datum follows a standard structure defined by
CIP-68:

\begin{lstlisting}[language={}, caption={CIP-68 datum structure (simplified)}]
{
  "constructor": 0,
  "fields": [
    {                          // metadata map
      "license_type": "PE",
      "jurisdiction": "US-CA",
      "license_number": "PE-2026-00042",
      "issue_date": "2026-01-15",
      "holder_pkh": "abc123...",
      "status": "active",      // <-- MUTABLE
      "last_updated": 12345678 // slot number
    },
    1                          // version number
  ]
}
\end{lstlisting}

The datum is encoded as CBOR (Concise Binary Object Representation) and
stored directly in the UTxO as an inline datum (CIP-32).  The \texttt{status}
field can hold any of: \texttt{active}, \texttt{suspended},
\texttt{revoked}, or \texttt{expired}.  Each status change creates a new
transaction on-chain, providing a complete audit trail.

The script address holding the Reference Token enforces that only the
authority's key can spend (and re-create) the UTxO.  This is typically
implemented as a simple \texttt{ScriptPubkey} validator or a Plutus V2
spending validator that checks for the authority's signature in the
transaction signatories.

\subsection{Verification Flow}

When a verifier wants to check a credential, the process is:

\begin{enumerate}[nosep]
  \item Observe the User Token (label~222) in the holder's wallet.
  \item Extract the policy ID and base asset name.
  \item Compute the Reference Token (label~100) asset name.
  \item Query the blockchain for the UTxO containing the Reference Token.
  \item Read the inline datum from that UTxO.
  \item Check the \texttt{status} field.
\end{enumerate}

This entire process requires only standard blockchain queries---no off-chain
service, no authority server, no API key.  The data is public, the datum is
inline (CIP-32), and the verification is purely mathematical.

\begin{figure}[ht]
\centering
\begin{tikzpicture}[
  entity/.style={guidebox, minimum width=3.2cm, minimum height=1.2cm},
  token/.style={guidebox, minimum width=3.8cm, minimum height=2.8cm,
    text width=3.4cm, align=left, font=\sffamily\scriptsize},
  node distance=1.2cm and 2.8cm,
]
  % --- Entities ---
  \node[entity, fill=purple!8] (authority) at (-3.5, 5.5)
    {\textbf{Licensing Board}\\(Authority)};
  \node[entity, fill=blue!8] (holder) at (3.5, 5.5)
    {\textbf{License Holder}\\(Professional)};

  % --- Reference Token ---
  \node[token, fill=orange!8] (reftok) at (-3.5, 1.8)
    {\textbf{Reference Token}\\(Label 100)\\[4pt]
    \texttt{policy: a1b2c3...}\\
    \texttt{name: (100)PE\_042}\\[3pt]
    \underline{Inline Datum:}\\
    \texttt{\ type: PE}\\
    \texttt{\ jurisdiction: US-CA}\\
    \texttt{\ status: \textcolor{green!50!black}{active}}\\
    \texttt{\ updated: slot\#789}};

  % --- User Token ---
  \node[token, fill=blue!8] (usertok) at (3.5, 1.8)
    {\textbf{User Token}\\(Label 222)\\[4pt]
    \texttt{policy: a1b2c3...}\\
    \texttt{name: (222)PE\_042}\\[3pt]
    \underline{CIP-25 Metadata:}\\
    \texttt{\ name: PE License}\\
    \texttt{\ image: ipfs://Qm..}\\[3pt]
    \footnotesize\textcolor{red!60!black}{IMMUTABLE}};

  % --- Script address ---
  \node[guidebox, fill=gray!10, minimum width=3.8cm, font=\sffamily\scriptsize]
    (script) at (-3.5, -0.8) {Script Address\\(authority-controlled)};

  % --- Personal wallet ---
  \node[guidebox, fill=gray!10, minimum width=3.8cm, font=\sffamily\scriptsize]
    (wallet) at (3.5, -0.8) {Personal Wallet\\(holder's private key)};

  % --- Arrows: ownership ---
  \draw[guidearrow] (authority) -- node[left, guidelabel, text width=1.8cm, align=center]
    {controls\\script} (reftok);
  \draw[guidearrow] (holder) -- node[right, guidelabel, text width=1.8cm, align=center]
    {holds in\\wallet} (usertok);
  \draw[guidearrow] (reftok) -- (script);
  \draw[guidearrow] (usertok) -- (wallet);

  % --- Update arrow ---
  \draw[guidearrow, color=red!60!black, very thick, dashed]
    (authority.south west) .. controls +(-1.5,-2) and +(-1.5,1.5) ..
    node[left, font=\sffamily\scriptsize, text=red!60!black, text width=2.2cm, align=center]
      {Can UPDATE\\datum to\\status: \textbf{revoked}}
    (reftok.west);

  % --- Link arrow ---
  \draw[guidearrow, color=green!50!black, <->]
    (reftok.east) -- node[above, guidelabel, text width=2.8cm, align=center]
      {Same policy ID\\+ base asset name\\= cryptographic link}
    (usertok.west);

  % --- Verifier ---
  \node[entity, fill=green!8] (verifier) at (0, -2.6)
    {\textbf{Verifier}\\(Any party)};
  \draw[guidearrow, dashed, color=green!50!black]
    (verifier.north) -- ++(0,0.7) -- ++(-3.5,0)
    node[above, near end, guidelabel] {reads datum}
    -- (script.south);
\end{tikzpicture}
\caption{CIP-68 dual-token pattern for credential management.  The authority
controls the Reference Token (label~100) at a script address and can update
its inline datum---including the status field---at any time.  The holder
possesses the immutable User Token (label~222) in their personal wallet.
A verifier reads the Reference Token's datum to determine current credential
status.  The authority never needs the holder's private key to revoke.}
\label{fig:cip68-dual-token}
\end{figure}

\subsection{CIP-25 vs.\ CIP-68: Comparison}

Table~\ref{tab:cip25-vs-cip68} summarizes the key differences between the
two metadata standards.  For the credential system, CIP-68's mutable datum
is the enabling feature that makes on-chain revocation practical.

\begin{table}[ht]
\centering
\caption{Comparison of CIP-25 and CIP-68 metadata standards}
\label{tab:cip25-vs-cip68}
\small
\begin{tabular}{@{}p{3.2cm}p{4.8cm}p{4.8cm}@{}}
\toprule
\textbf{Property} & \textbf{CIP-25} & \textbf{CIP-68} \\
\midrule
Metadata location & Transaction auxiliary data (label~721) & Inline datum on Reference Token UTxO \\
Mutability & \textbf{Immutable} --- written once at mint & \textbf{Mutable} --- authority can update datum \\
Token count & 1 token per credential & 2 tokens: Reference (100) + User (222) \\
Who holds metadata & Embedded in minting tx (no custodian) & Reference Token at authority script address \\
Revocation & Not possible on-chain & Datum update: \texttt{status $\to$ revoked} \\
Authority key needed & Only at mint time & At mint \emph{and} for every datum update \\
Holder key needed & Only to transfer/burn the token & Only to transfer/burn the User Token \\
On-chain cost & $\sim$0.2~ADA (single token mint) & $\sim$0.4~ADA (two tokens + datum) \\
Reading metadata & Parse minting tx auxiliary data & Read inline datum from Reference UTxO \\
Audit trail & Single minting event & Full history of datum updates on-chain \\
Requires CIP-31 & No & Yes (reference inputs for reading datum) \\
Requires CIP-32 & No & Yes (inline datums for storing datum) \\
\bottomrule
\end{tabular}
\end{table}

\subsection{Why CIP-68 Matters for Credentials}

The importance of CIP-68 for the credential system cannot be overstated.
Before CIP-68, any Cardano-based credential project faced an impossible
trilemma: (1)~you could store metadata on-chain but it was immutable; (2)~you
could store metadata off-chain but that reintroduced centralization; or
(3)~you could accept that revocation was not truly decentralized.

CIP-68 breaks this trilemma by enabling \textbf{mutable on-chain metadata
under authority control}.  The result is a system where:

\begin{itemize}[nosep]
  \item Credential issuance is on-chain and cryptographically verifiable.
  \item Credential status is on-chain and always current.
  \item Revocation is unilateral by the authority---no holder cooperation
    required.
  \item Verification is fully decentralized---any node can read the datum.
  \item The audit trail is immutable---every status change is a new
    transaction.
\end{itemize}

The dual-token pattern also provides a natural separation between
\emph{proof of issuance} (the User Token in the holder's wallet says ``this
authority issued me a credential'') and \emph{current status} (the Reference
Token's datum says ``the credential is currently active/revoked/suspended'').
This separation mirrors real-world credential systems where a paper
certificate proves issuance but a database query confirms current status---
except that in the CIP-68 model, both proofs are on-chain.


% ============================================================================
\section{CIP-30: dApp--Wallet Web Bridge}
\label{sec:cip30}
% ============================================================================

\subsection{The Problem}

A web-based credential management application needs to interact with a
user's Cardano wallet: reading their address, checking token balances,
and---most critically---asking the wallet to sign transactions.  Without a
standard API, every wallet (Eternl, Nami, Lace, Flint, Typhon, etc.) would
expose its own proprietary JavaScript interface, forcing dApp developers to
write wallet-specific code and creating fragmentation across the ecosystem.

More importantly, the security model must be \textbf{non-custodial}: the
dApp must \emph{never} see the user's private keys.  If the application
server were compromised, an attacker with access to private keys could drain
wallets, forge signatures, or transfer credentials.  The wallet must sign
transactions locally and return only the signed transaction to the dApp.

\subsection{How It Works}

CIP-30 defines a JavaScript API that Cardano wallets inject into the
browser's \texttt{window.cardano} namespace.  The API follows a two-phase
pattern:

\begin{enumerate}[nosep]
  \item \textbf{Discovery and Connection.}  The dApp checks
    \texttt{window.cardano} for available wallets, then calls
    \texttt{cardano.<wallet>.enable()} to request user authorization.  The
    user sees a popup in their wallet extension asking whether to connect.
    If approved, the wallet returns an API handle.
  \item \textbf{Interaction.}  The dApp uses the API handle to call
    methods such as \texttt{getUsedAddresses()},
    \texttt{getUtxos()}, \texttt{getBalance()}, and---most
    importantly---\texttt{signTx(tx, partialSign)}.  The
    \texttt{signTx} method takes an unsigned transaction (CBOR-encoded),
    displays it to the user for approval, signs it with the wallet's private
    key, and returns the signature (or signed transaction).
\end{enumerate}

The critical security property is that the private key \emph{never leaves
the wallet extension}.  The dApp constructs the transaction, but the wallet
signs it.  This is analogous to how a hardware wallet works: the application
builds the transaction, the device signs it, and the application never sees
the signing key.

\subsection{How the Credential System Uses CIP-30}

In the credential system, licensees and signers connect their wallets via
CIP-30 to the web dashboard.  When a professional needs to sign a document
(depositing signature and validity tokens into the work product validator),
the application:

\begin{enumerate}[nosep]
  \item Builds the transaction off-chain using PyCardano (or a JavaScript
    library in the browser).
  \item Sends the unsigned transaction CBOR to the wallet via
    \texttt{signTx()}.
  \item The wallet displays the transaction details (tokens being sent,
    destination address, fee) and asks the user to confirm.
  \item The user approves; the wallet signs and returns the witness.
  \item The application attaches the witness and submits the transaction to
    the network.
\end{enumerate}

At no point does the application server have access to the user's signing
key.  Even if the server is compromised, the attacker cannot sign
transactions on behalf of users.  This non-custodial model is essential for
a credential system where users must maintain sovereign control over their
professional identities.

\begin{figure}[ht]
\centering
\begin{tikzpicture}[
  actor/.style={guidebox, minimum width=3cm, minimum height=1.2cm},
  node distance=1.5cm and 4.5cm,
]
  % Actors
  \node[actor, fill=blue!8] (dapp) at (0,0) {\textbf{dApp}\\(Web Browser)};
  \node[actor, fill=green!8] (wallet) at (6,0) {\textbf{Wallet Extension}\\(Eternl / Nami / Lace)};
  \node[actor, fill=purple!8] (chain) at (3,-5.5) {\textbf{Cardano Network}};

  % Lifelines
  \draw[dashed, gray!40] (dapp.south) -- ++(0,-5);
  \draw[dashed, gray!40] (wallet.south) -- ++(0,-5);

  % Interactions
  \draw[guidearrow] (0.1,-0.8) -- node[above, guidelabel]
    {1. \texttt{cardano.eternl.enable()}} (5.9,-0.8);
  \draw[guidearrow] (5.9,-1.4) -- node[above, guidelabel]
    {2. API handle (user approved)} (0.1,-1.4);
  \draw[guidearrow] (0.1,-2.0) -- node[above, guidelabel]
    {3. \texttt{api.getUsedAddresses()}} (5.9,-2.0);
  \draw[guidearrow] (5.9,-2.6) -- node[above, guidelabel]
    {4. \texttt{[addr\_test1q...]}} (0.1,-2.6);
  \draw[guidearrow] (0.1,-3.2) -- node[above, guidelabel]
    {5. \texttt{api.signTx(unsignedTxCBOR)}} (5.9,-3.2);
  \draw[guidearrow] (5.9,-3.8) -- node[above, guidelabel]
    {6. Signed witness (key \textbf{never} leaves wallet)} (0.1,-3.8);
  \draw[guidearrow] (0.1,-4.5) -- node[below left, guidelabel, text width=2.5cm, align=center]
    {7. Submit\\signed tx} (chain);

  % Security note
  \node[draw, rounded corners=2pt, fill=red!6, font=\sffamily\scriptsize,
    text width=4.5cm, align=center] at (3,-7)
    {Private key NEVER leaves the wallet.\\The dApp sees only the signature.};
\end{tikzpicture}
\caption{CIP-30 dApp--Wallet interaction flow.  The web application requests
connection, reads addresses and UTxOs, and sends unsigned transactions to the
wallet for signing.  The wallet extension holds the private key locally and
returns only the cryptographic witness.  The dApp never has access to the
signing key.}
\label{fig:cip30-flow}
\end{figure}


% ============================================================================
\section{CIP-31: Reference Inputs}
\label{sec:cip31}
% ============================================================================

\subsection{The Problem}

In the original eUTxO model, every UTxO referenced by a transaction is
\textbf{consumed} (spent).  If a transaction needs to read data from a UTxO
---for example, checking a CIP-68 Reference Token's datum to verify a
credential's status---it must spend that UTxO and recreate it in the same
transaction.  This creates three problems:

\begin{enumerate}[nosep]
  \item \textbf{Concurrency.}  If two transactions both need to read the
    same Reference Token datum, they both try to spend the same UTxO.  Only
    one can succeed; the other fails with a UTxO conflict.  For a credential
    that many parties may verify simultaneously, this is unacceptable.
  \item \textbf{Complexity.}  The spending transaction must include the
    script that guards the UTxO, a redeemer, and must recreate the UTxO with
    exactly the same value and datum---adding transaction size, fees, and
    potential for bugs.
  \item \textbf{Authorization.}  Spending a UTxO requires satisfying its
    spending conditions (e.g., the authority's signature).  A verifier who
    merely wants to \emph{read} the datum should not need the authority's
    key.
\end{enumerate}

\subsection{How It Works}

CIP-31 introduces \textbf{reference inputs}: a new transaction field that
allows a transaction to \emph{read} a UTxO without consuming it.  The UTxO
remains untouched on the ledger after the transaction is processed.  The
transaction's scripts can access the reference input's value, datum, and
address, but the UTxO is not spent.

This is a Plutus~V2 feature introduced in the Vasil hard fork.  In the
transaction body, reference inputs are listed in a separate field
(\texttt{reference\_inputs}) distinct from regular inputs (\texttt{inputs}).
The ledger validates that each referenced UTxO exists at the time the
transaction is processed but does not remove it from the UTxO set.

\subsection{How the Credential System Uses CIP-31}

CIP-31 is critical for the credential system's verification flow.  When a
verifier or a smart contract needs to check a license's current status, the
transaction includes the CIP-68 Reference Token UTxO as a \textbf{reference
input}.  The script reads the inline datum, checks the \texttt{status} field,
and proceeds---all without spending the Reference Token.

This means:
\begin{itemize}[nosep]
  \item Multiple verifications can happen simultaneously (no UTxO contention).
  \item The verifier does not need the authority's key.
  \item The Reference Token UTxO remains stable on-chain.
  \item Transaction size is smaller (no need to recreate the UTxO).
\end{itemize}

Without CIP-31, the CIP-68 pattern would be impractical for any credential
that is verified frequently, because each verification would compete to spend
the same Reference Token UTxO.

\begin{figure}[ht]
\centering
\begin{tikzpicture}[
  utxo/.style={guidebox, minimum width=3.4cm, minimum height=1.4cm,
    font=\sffamily\scriptsize},
  tx/.style={guidebox, fill=yellow!8, minimum width=2.8cm, minimum height=1.2cm},
  node distance=0.8cm and 2.2cm,
]
  % --- Consuming (old way) ---
  \node[font=\sffamily\small\bfseries] at (-4.5,4.2) {Consuming Input (Pre-CIP-31)};

  \node[utxo, fill=orange!8] (utxo1) at (-6,2.8) {Reference Token\\UTxO\\datum: status=active};
  \node[tx] (tx1) at (-3,2.8) {Transaction};
  \node[utxo, fill=orange!8] (utxo1b) at (-3,0.8) {Recreated UTxO\\(same datum)};

  \draw[guidearrow] (utxo1) -- node[above, guidelabel] {consumed} (tx1);
  \draw[guidearrow, dashed] (tx1) -- node[right, guidelabel] {must recreate} (utxo1b);
  \draw[-{Latex[length=2.5mm]}, very thick, red!60!black]
    (-6.8,2.1) -- (-6.8,1.5) node[below, font=\sffamily\tiny, text=red!60!black]
    {UTxO destroyed};

  % --- Reference Input (CIP-31) ---
  \node[font=\sffamily\small\bfseries] at (3.5,4.2) {Reference Input (CIP-31)};

  \node[utxo, fill=green!8] (utxo2) at (2,2.8) {Reference Token\\UTxO\\datum: status=active};
  \node[tx] (tx2) at (5,2.8) {Transaction};

  \draw[guidearrow, dashed, color=green!50!black] (utxo2) --
    node[above, guidelabel] {read only} (tx2);
  \draw[-{Latex[length=2.5mm]}, very thick, green!50!black]
    (1.2,2.1) -- (1.2,1.5) node[below, font=\sffamily\tiny, text=green!50!black]
    {UTxO survives};

  \node[utxo, fill=green!8] (utxo2b) at (2,0.8) {Same UTxO\\still on ledger\\(untouched)};

  % Divider
  \draw[dashed, gray!40] (-0.5,4.5) -- (-0.5,-0.2);
\end{tikzpicture}
\caption{CIP-31 reference inputs compared to consuming inputs.  On the left,
the pre-CIP-31 approach: reading a datum requires consuming the UTxO and
recreating it (expensive, causes contention).  On the right, CIP-31: the
transaction reads the UTxO without consuming it---the UTxO survives
untouched, enabling concurrent reads by multiple transactions.}
\label{fig:cip31-reference-inputs}
\end{figure}


% ============================================================================
\section{CIP-32: Inline Datums}
\label{sec:cip32}
% ============================================================================

\subsection{The Problem}

Before CIP-32, the eUTxO model supported attaching data to script-locked
UTxOs through \textbf{datum hashes}.  When creating a UTxO at a script
address, the transaction included only the \emph{hash} of the datum in the
UTxO itself.  The actual datum content had to be provided separately: either
embedded elsewhere in the transaction's witness set, or stored entirely
off-chain and supplied by whoever eventually spent the UTxO.

This design had two major consequences:

\begin{enumerate}[nosep]
  \item \textbf{Off-chain burden.}  The party spending the UTxO had to know
    the original datum and include it in their spending transaction.  This
    required off-chain coordination: the datum had to be stored somewhere
    (a database, IPFS, a centralized service) and communicated to the
    spender.  If the datum was lost, the UTxO became unspendable.
  \item \textbf{No public readability.}  A third party observing the chain
    could see that a UTxO had a datum hash, but could not determine the
    datum's contents without access to the original data.  This made on-chain
    data effectively opaque to verifiers who had not participated in the
    original transaction.
\end{enumerate}

For a credential system where \emph{any party} must be able to read a
credential's current status, datum hashes are fundamentally inadequate.  A
verifier cannot check \texttt{status: active} if they can only see
\texttt{datum\_hash: 0x3f7a...}.

\subsection{How It Works}

CIP-32 introduces \textbf{inline datums}: the full datum content is stored
\emph{directly in the UTxO} on the ledger.  Any node, wallet, or explorer
that can read the UTxO can also read the datum---no off-chain lookup
required, no special access, no coordination with the datum creator.

Technically, a UTxO's datum field now supports three states:
\begin{description}[nosep, leftmargin=1.5em]
  \item[None] No datum attached (standard ADA-only UTxO).
  \item[Datum Hash] Only the hash is stored on-chain (pre-CIP-32 behavior).
  \item[Inline Datum] The full datum is stored on-chain (CIP-32).
\end{description}

The trade-off is that inline datums increase the UTxO's size on the ledger
and therefore increase the minimum ADA required (minUTxO) to hold the UTxO.
For a credential datum containing license type, jurisdiction, status, and a
few additional fields, the size increase is modest (a few hundred bytes),
and the minUTxO impact is approximately 0.2--0.5~ADA additional.

\subsection{How the Credential System Uses CIP-32}

CIP-32 is what makes CIP-68 \emph{possible}.  The CIP-68 Reference Token
stores its metadata as an inline datum so that:

\begin{itemize}[nosep]
  \item Any verifier can read the credential status by querying the UTxO.
  \item No off-chain service is needed to resolve datum hashes.
  \item Smart contracts executing on-chain (e.g., the Signature Collection
    Validator) can read the Reference Token's datum via reference inputs
    (CIP-31) and make validation decisions based on the credential's current
    status.
\end{itemize}

Without inline datums, the dual-token pattern would require an off-chain
datum registry---reintroducing exactly the centralized infrastructure that
the blockchain architecture is designed to eliminate.  CIP-32, together with
CIP-31, forms the technical foundation upon which CIP-68's mutable on-chain
metadata is built.


% ============================================================================
\section{CIP-33: Reference Scripts}
\label{sec:cip33}
% ============================================================================

\subsection{The Problem}

Plutus scripts (smart contracts) can be large.  A compiled Plutus~V2 script
may be 5--15~KB of serialized CBOR.  Before CIP-33, every transaction that
invoked a Plutus script had to include the \emph{entire compiled script} in
its witness set.  For a credential system where every signing operation
invokes the Signature Collection Validator, this means attaching the same
multi-kilobyte script to every transaction---wasting block space, increasing
fees, and potentially pushing transactions over the 16,384-byte maximum
transaction size.

\subsection{How It Works}

CIP-33 allows a Plutus script to be \textbf{deployed once to a UTxO} and
then \textbf{referenced} by subsequent transactions.  The deployment
transaction creates a UTxO containing the script in a special
\texttt{reference\_script} field.  Future transactions that need to execute
the script include only a reference to this UTxO (its transaction hash and
output index), not the script itself.

The ledger resolves the reference at validation time: when processing a
transaction that references a script UTxO, the node fetches the script from
the referenced UTxO and uses it for validation.  The referenced UTxO is
\emph{not consumed}---like CIP-31 reference inputs, the script UTxO
persists on-chain and can be referenced by unlimited future transactions.

\subsection{How the Credential System Uses CIP-33}

The credential system deploys the Signature Collection Validator and the
Dues Enforcement Contract as reference scripts.  This provides three
benefits:

\begin{enumerate}[nosep]
  \item \textbf{Smaller transactions.}  Signing transactions include only a
    32-byte reference (transaction hash) plus an index, instead of a
    multi-kilobyte script.  This can reduce transaction size by 80--90\%.
  \item \textbf{Lower fees.}  Cardano transaction fees are proportional to
    transaction size.  Smaller transactions mean lower fees for every
    signing operation.
  \item \textbf{Auditability.}  The script is deployed once at a known UTxO.
    Anyone can inspect the exact validator code that governs the credential
    system by reading this single UTxO.
\end{enumerate}

\begin{figure}[ht]
\centering
\begin{tikzpicture}[
  txbox/.style={guidebox, minimum width=4cm, minimum height=2.6cm,
    text width=3.6cm, align=left, font=\sffamily\scriptsize},
  scriptbox/.style={guidebox, fill=purple!8, minimum width=3.2cm,
    minimum height=1.2cm, font=\sffamily\scriptsize},
]
  % --- Without CIP-33 ---
  \node[font=\sffamily\small\bfseries] at (-4.5,4.2)
    {Without CIP-33};

  \node[txbox, fill=red!6] (tx1) at (-4.5,1.6)
    {\textbf{Transaction \#1}\\[2pt]
     inputs: [...]\\
     outputs: [...]\\
     \textcolor{red!60!black}{witness: \{full\_script}\\
     \textcolor{red!60!black}{\ \ (8 KB of CBOR)\}}\\[2pt]
     size: $\sim$9.5 KB};

  \node[txbox, fill=red!6] (tx2) at (-4.5,-2)
    {\textbf{Transaction \#2}\\[2pt]
     inputs: [...]\\
     outputs: [...]\\
     \textcolor{red!60!black}{witness: \{full\_script}\\
     \textcolor{red!60!black}{\ \ (8 KB of CBOR)\}}\\[2pt]
     size: $\sim$9.5 KB};

  % --- With CIP-33 ---
  \node[font=\sffamily\small\bfseries] at (3.5,4.2)
    {With CIP-33};

  \node[scriptbox] (deploy) at (3.5,3)
    {Script UTxO\\(deployed once)\\8 KB};

  \node[txbox, fill=green!8] (tx3) at (3.5,0.6)
    {\textbf{Transaction \#1}\\[2pt]
     inputs: [...]\\
     outputs: [...]\\
     \textcolor{green!50!black}{ref\_script: tx\#abc.0}\\
     \textcolor{green!50!black}{\ \ (32-byte ref)}\\[2pt]
     size: $\sim$1.5 KB};

  \node[txbox, fill=green!8] (tx4) at (3.5,-2.6)
    {\textbf{Transaction \#2}\\[2pt]
     inputs: [...]\\
     outputs: [...]\\
     \textcolor{green!50!black}{ref\_script: tx\#abc.0}\\
     \textcolor{green!50!black}{\ \ (32-byte ref)}\\[2pt]
     size: $\sim$1.5 KB};

  % Arrows
  \draw[guidearrow, dashed, color=green!50!black]
    (tx3.north) -- ++(0,0.5) -| (deploy.south);
  \draw[guidearrow, dashed, color=green!50!black]
    (tx4.north east) -- ++(0.5,0.5) -- ++(0,3.5) -| (deploy.east);

  % Savings annotation
  \node[draw, rounded corners=2pt, fill=yellow!8, font=\sffamily\scriptsize,
    text width=3cm, align=center] at (3.5,-4.5)
    {Savings: $\sim$80--90\% smaller\\transactions after deployment};

  % Divider
  \draw[dashed, gray!40] (-0.5,4.5) -- (-0.5,-5);
\end{tikzpicture}
\caption{CIP-33 reference scripts.  Without CIP-33 (left), every transaction
must include the full compiled Plutus script in its witness set, consuming
kilobytes of space per transaction.  With CIP-33 (right), the script is
deployed once to a UTxO, and subsequent transactions reference it with a
compact 32-byte pointer.  Transaction size drops by 80--90\%.}
\label{fig:cip33-reference-scripts}
\end{figure}


% ============================================================================
\section{CIP-1852 / BIP-44: HD Wallet Derivation}
\label{sec:cip1852}
% ============================================================================

\subsection{The Problem}

A naive approach to key management would generate a single random private
key and derive a single address from it.  This creates several problems:

\begin{enumerate}[nosep]
  \item \textbf{No backup recoverability.}  If the key file is lost, all
    funds and credentials are permanently inaccessible.  There is no way to
    regenerate the key.
  \item \textbf{Privacy.}  Using a single address for all transactions
    creates a complete, publicly visible transaction graph.  Every
    credential, every payment, every interaction is linked.
  \item \textbf{Key management complexity.}  If a user needs separate keys
    for different purposes (payment, staking, multiple credentials), they
    must manage multiple independent keys with independent backups.
  \item \textbf{No standard.}  Without a derivation convention, every wallet
    would generate keys differently, making wallet recovery across different
    software impossible.
\end{enumerate}

\subsection{How It Works}

CIP-1852 defines Cardano's \textbf{hierarchical deterministic (HD) wallet}
derivation scheme, building on the BIP-32 (key derivation), BIP-39
(mnemonic encoding), and BIP-44 (multi-account structure) standards from the
Bitcoin ecosystem.

The core idea is that a single \textbf{24-word mnemonic phrase} (BIP-39)
encodes a master seed from which an unlimited number of keys can be derived
deterministically.  Given the same mnemonic, any compliant wallet software
will derive exactly the same keys in exactly the same order.  This means:

\begin{itemize}[nosep]
  \item \textbf{Backup is trivial.}  Write down 24 words.  That is the
    complete backup for all current and future keys.
  \item \textbf{Recovery works across wallets.}  Import the mnemonic into
    Eternl, Nami, Lace, or any CIP-1852-compliant wallet, and all keys are
    regenerated identically.
  \item \textbf{Unlimited addresses.}  Increment the index in the derivation
    path to generate fresh addresses for privacy.
\end{itemize}

\subsubsection{The Derivation Path}

CIP-1852 defines the following derivation path structure:

\begin{center}
\texttt{m / purpose' / coin\_type' / account' / role / index}
\end{center}

\begin{description}[nosep, leftmargin=1.5em]
  \item[\texttt{m}] The master key derived from the BIP-39 mnemonic.
  \item[\texttt{purpose' = 1852'}] Identifies this as a Cardano Shelley-era
    derivation (named after Ada Lovelace's birth year).  The apostrophe
    denotes \emph{hardened derivation}.
  \item[\texttt{coin\_type' = 1815'}] Identifies the Cardano blockchain
    (named after Ada Lovelace's death year).  Registered in the SLIP-44
    coin type registry.
  \item[\texttt{account' = 0'}] Account index.  Most users use account~0;
    multiple accounts allow logical separation (e.g., personal vs.\
    professional credentials).
  \item[\texttt{role}] Distinguishes key purpose:
    \begin{itemize}[nosep]
      \item \texttt{0} = External payment keys (receiving addresses)
      \item \texttt{1} = Internal payment keys (change addresses)
      \item \texttt{2} = Staking keys
    \end{itemize}
  \item[\texttt{index}] Sequential index within the role.  Incrementing
    this generates additional keys.
\end{description}

The two most important paths for the credential system are:

\begin{center}
\begin{tabular}{ll}
  \texttt{m/1852'/1815'/0'/0/0} & First payment key (for receiving tokens) \\
  \texttt{m/1852'/1815'/0'/2/0} & First staking key (for delegation) \\
\end{tabular}
\end{center}

\subsubsection{Hardened vs.\ Unhardened Derivation}

The apostrophe (\texttt{'}) in the derivation path denotes \textbf{hardened
derivation}.  This is a critical security distinction:

\begin{description}[nosep, leftmargin=1.5em]
  \item[Unhardened (normal) derivation] uses the parent's \emph{public key}
    to derive child keys.  This means anyone with the parent public key can
    compute all child public keys.  Useful for generating receiving addresses
    without exposing the private key (e.g., a watch-only wallet or a payment
    server).
  \item[Hardened derivation] uses the parent's \emph{private key} to derive
    child keys.  Knowing the parent public key reveals \emph{nothing} about
    child keys.  This provides a security firewall: if a child private key
    is compromised, the attacker cannot compute sibling or parent keys.
\end{description}

CIP-1852 requires hardened derivation for purpose, coin type, and account
levels.  This ensures that compromise of a single payment key does not
allow an attacker to derive the staking key, other account keys, or the
master key.  The role and index levels use unhardened derivation, which
enables watch-only wallets to generate addresses for receiving payments
without access to private keys.

\subsection{How the Credential System Uses CIP-1852}

Every wallet in the credential system---authority, licensee, signer, and
observer---is generated from a 24-word BIP-39 mnemonic using CIP-1852
derivation.  The \texttt{generate\_wallet()} function:

\begin{enumerate}[nosep]
  \item Generates or accepts a BIP-39 mnemonic (24 words, 256 bits of
    entropy).
  \item Derives the master key using PBKDF2-HMAC-SHA512.
  \item Derives the payment key at \texttt{m/1852'/1815'/0'/0/0}.
  \item Derives the staking key at \texttt{m/1852'/1815'/0'/2/0}.
  \item Constructs the base address by combining the payment key hash (PKH,
    28~bytes) and staking key hash (SKH, 28~bytes) with the network
    identifier.
  \item Encodes the address in bech32 format (\texttt{addr\_test1...} for
    testnet, \texttt{addr1...} for mainnet).
\end{enumerate}

The mnemonic is the \textbf{single point of backup}: losing it means losing
all keys and all credentials.  In the production deployment, authority wallet
mnemonics are encrypted with AES-256-GCM and stored with strict access
controls.  Licensee and signer mnemonics are managed by their personal
wallets (Eternl, Nami, etc.) via CIP-30---the credential system never sees
or stores these mnemonics.

\begin{figure}[ht]
\centering
\begin{tikzpicture}[
  level 1/.style={sibling distance=5.5cm, level distance=1.6cm},
  level 2/.style={sibling distance=4cm, level distance=1.6cm},
  level 3/.style={sibling distance=5cm, level distance=1.6cm},
  level 4/.style={sibling distance=2.8cm, level distance=1.6cm},
  level 5/.style={sibling distance=1.6cm, level distance=1.6cm},
  treenode/.style={guidebox, minimum width=1.8cm, minimum height=0.7cm,
    font=\sffamily\tiny, inner sep=2pt},
  edge from parent/.style={guidearrow, thick},
  edge from parent path={(\tikzparentnode.south) -- (\tikzchildnode.north)},
]
  \node[treenode, fill=purple!8, minimum width=4cm]
    {\textbf{Master Key (m)}\\24-word mnemonic $\to$ seed}
    child {
      node[treenode, fill=blue!8, minimum width=2.8cm]
        {\textbf{1852'} (purpose)\\Cardano Shelley\\{\tiny\textcolor{red!60!black}{hardened}}}
      child {
        node[treenode, fill=blue!8, minimum width=2.5cm]
          {\textbf{1815'} (coin type)\\Cardano ADA\\{\tiny\textcolor{red!60!black}{hardened}}}
        child {
          node[treenode, fill=blue!8, minimum width=2.4cm]
            {\textbf{0'} (account)\\Default account\\{\tiny\textcolor{red!60!black}{hardened}}}
          child {
            node[treenode, fill=green!8, minimum width=2cm]
              {\textbf{0} (external)\\Payment keys}
            child {
              node[treenode, fill=orange!8, minimum width=1.6cm]
                {\textbf{0}\\1st payment\\key}
            }
            child {
              node[treenode, fill=orange!8, minimum width=1.6cm]
                {\textbf{1}\\2nd payment\\key}
            }
          }
          child {
            node[treenode, fill=green!8, minimum width=2cm]
              {\textbf{2} (staking)\\Stake keys}
            child {
              node[treenode, fill=orange!8, minimum width=1.6cm]
                {\textbf{0}\\1st stake\\key}
            }
          }
        }
      }
    };

  % Annotations
  \node[guidelabel, text width=2.5cm, align=left] at (6.5,-2.4)
    {Hardened levels\\(marked with ')\\use private key\\for derivation};
  \node[guidelabel, text width=2.5cm, align=left] at (6.5,-5.6)
    {Unhardened levels\\use public key\\$\to$ watch-only\\wallets possible};

  % Bracket for hardened
  \draw[decorate, decoration={brace, amplitude=5pt, mirror},
    thick, gray!50!black]
    (5.5,-0.5) -- (5.5,-4.2) node[midway, right=8pt] {};

  % Bracket for unhardened
  \draw[decorate, decoration={brace, amplitude=5pt, mirror},
    thick, gray!50!black]
    (5.5,-4.8) -- (5.5,-7.6) node[midway, right=8pt] {};
\end{tikzpicture}
\caption{CIP-1852 / BIP-44 hierarchical deterministic wallet derivation tree.
The master key is derived from a 24-word mnemonic.  Purpose (1852'),
coin type (1815'), and account (0') levels use hardened derivation (private
key required), providing a security firewall.  Role and index levels use
unhardened derivation, enabling watch-only wallets.  The credential system
uses \texttt{m/1852'/1815'/0'/0/0} for payment and
\texttt{m/1852'/1815'/0'/2/0} for staking.}
\label{fig:cip1852-derivation}
\end{figure}


% ============================================================================
\section{How the CIPs Work Together}
\label{sec:cips-together}
% ============================================================================

The CIPs described in this chapter are not isolated features---they form an
interlocking system where each standard enables or enhances the others.
Understanding their interactions is essential for understanding the
credential architecture as a whole.

\textbf{CIP-32 enables CIP-68.}  Without inline datums, the Reference
Token's metadata would be stored as a hash, requiring off-chain resolution.
Inline datums make the metadata directly readable on-chain.

\textbf{CIP-31 makes CIP-68 practical.}  Without reference inputs, every
credential verification would consume (and have to recreate) the Reference
Token UTxO, causing contention.  Reference inputs allow unlimited concurrent
reads.

\textbf{CIP-33 makes CIP-68 affordable.}  Without reference scripts, every
transaction interacting with the credential validator would include the full
Plutus script, inflating transaction size and fees.  Reference scripts reduce
per-transaction overhead to a 32-byte pointer.

\textbf{CIP-68 makes CIP-25 revocable.}  CIP-25 provides immutable
attestation (``this credential was issued'').  CIP-68 adds mutable status
(``the credential is currently active/revoked'').  Together, they provide
both historical provenance and current validity.

\textbf{CIP-30 connects users to the system.}  Without the dApp--wallet
bridge, every credential operation would require server-side key management,
creating custodial risk.  CIP-30 keeps private keys under user control.

\textbf{CIP-1852 makes wallets recoverable and interoperable.}  Without
HD derivation, losing a key file means losing all credentials.  With
CIP-1852, 24 words are sufficient to recover everything.

\begin{figure}[ht]
\centering
\begin{tikzpicture}[
  cipbox/.style={guidebox, minimum width=2.8cm, minimum height=1cm,
    font=\sffamily\scriptsize},
  node distance=0.6cm and 1.2cm,
]
  % Central: CIP-68
  \node[cipbox, fill=orange!8, minimum width=3.2cm, minimum height=1.4cm,
    font=\sffamily\small] (cip68) at (0,0)
    {\textbf{CIP-68}\\Mutable Metadata};

  % Dependencies feeding into CIP-68
  \node[cipbox, fill=green!8] (cip32) at (-4, 1.5) {CIP-32\\Inline Datums};
  \node[cipbox, fill=green!8] (cip31) at (-4,-1.5) {CIP-31\\Reference Inputs};
  \node[cipbox, fill=green!8] (cip33) at (-4, 0) {CIP-33\\Reference Scripts};

  % CIP-25 companion
  \node[cipbox, fill=blue!8] (cip25) at (4, 1.5) {CIP-25\\Immutable Metadata};

  % CIP-30 and CIP-1852
  \node[cipbox, fill=purple!8] (cip30) at (4, 0) {CIP-30\\Wallet Bridge};
  \node[cipbox, fill=yellow!8] (cip1852) at (4,-1.5) {CIP-1852\\HD Wallets};

  % Credential System
  \node[cipbox, fill=red!6, minimum width=4cm, minimum height=1.2cm,
    font=\sffamily\small] (system) at (0,-3.5)
    {\textbf{Credential Verification System}};

  % Arrows: dependencies
  \draw[guidearrow] (cip32) -- node[above, guidelabel] {stores datum} (cip68);
  \draw[guidearrow] (cip31) -- node[below, guidelabel] {reads datum} (cip68);
  \draw[guidearrow] (cip33) -- node[above, guidelabel] {deploys script} (cip68);

  % Arrows: to system
  \draw[guidearrow] (cip68) -- (system);
  \draw[guidearrow] (cip25) |- node[near end, above, guidelabel] {attestation} (system.north east);
  \draw[guidearrow] (cip30) -- (system);
  \draw[guidearrow] (cip1852) |- (system.south east);
\end{tikzpicture}
\caption{CIP dependency graph for the credential verification system.
CIPs~31, 32, and~33 (the Plutus~V2 trio) form the technical foundation that
makes CIP-68's mutable datum pattern practical and affordable.  CIP-25
provides immutable attestation.  CIP-30 and CIP-1852 handle wallet
connectivity and key management.  All six converge in the credential
verification system.}
\label{fig:cip-dependency}
\end{figure}


% ============================================================================
\section{Chapter Summary}
\label{sec:cips-summary}
% ============================================================================

This chapter examined the eight CIPs that form the technical foundation of
the credential verification system:

\begin{enumerate}[nosep]
  \item \textbf{CIP-25} provides immutable NFT metadata, establishing
    permanent on-chain attestation of credential issuance.
  \item \textbf{CIP-68} introduces the dual-token pattern with mutable
    inline datums, enabling authority-controlled revocation without requiring
    the holder's private key---the key innovation for on-chain credential
    management.
  \item \textbf{CIP-30} defines the non-custodial dApp--wallet bridge,
    keeping private keys under user control while enabling web-based
    credential operations.
  \item \textbf{CIP-31} introduces reference inputs, allowing transactions
    to read UTxO data without consuming it---essential for concurrent
    credential verification.
  \item \textbf{CIP-32} introduces inline datums, storing full data directly
    in UTxOs---the prerequisite for CIP-68's on-chain metadata.
  \item \textbf{CIP-33} introduces reference scripts, reducing transaction
    size and fees by deploying Plutus scripts once and referencing them.
  \item \textbf{CIP-1852} defines HD wallet derivation, providing
    recoverable, interoperable key management from a single mnemonic.
\end{enumerate}

The next chapter will examine how these CIPs are combined in the system
architecture: the smart contract layer, the token taxonomy, and the
transaction workflows that tie everything together.
